% !TEX TS-program = pdflatex
% !TEX encoding = UTF-8 Unicode
% !BIB TS-program = bibtex
%
% ============================================================================
% paper.tex -- {{paper-title}}
%
% Target  : {{journal}} (IOP Publishing)
% Class   : iopjournal.cls (2024/01/31), supplied by IOP at
%   https://publishingsupport.iopscience.iop.org/questions/latex-template/
% Authors : {{authors}}, CESI LINEACT, {{year}}.
%
% Companion files :
%   paper_si.tex                 -- Supplementary Information
%   references/references.bib    -- BibTeX
%   figures/                     -- publication figures (PDF)
%   experiments/d{NN}_{name}.py  -- reproducible scripts
%
% Compile :
%   pdflatex paper.tex
%   bibtex   paper
%   pdflatex paper.tex
%   pdflatex paper.tex
% ============================================================================
\documentclass{iopjournal}      % add [anonymous] for double-anonymous review

\usepackage[T1]{fontenc}
\usepackage[utf8]{inputenc}
\usepackage{microtype}
\usepackage{csquotes}
\usepackage{amssymb,amsmath,amstext,amsthm}

\graphicspath{{figures/}}

\usepackage{booktabs}
\usepackage{array}
\usepackage{float}
\usepackage{enumitem}
\setlist{topsep=2pt,itemsep=2pt,parsep=0pt}
\usepackage{url}

% IOP accepts either Vancouver (numbered) or Harvard (author-year).
% Defaulting to author-year + plainnat ; switch to iopart-num for numeric.
\usepackage[authoryear,round]{natbib}
\bibliographystyle{plainnat}

% Patch \orcid to a text fallback if the orcid.pdf image is missing.
\IfFileExists{orcid.pdf}{}{%
  \renewcommand{\orcid}[1]{\textsuperscript{\href{https://orcid.org/##1}{\textsf{[ORCID]}}}}%
}

% Running-header customisation : IOP class hardcodes "Journal vv (yyyy) aaaaaa"
% and "Author et al" ; override for the target journal.
\usepackage{fancyhdr}
\fancyhead[L]{{\small \sf IOP Publishing}\hspace{5mm}%
              {\it {{journal-short}}}\ \textbf{vv} (yyyy) aaaaaa}
\fancyhead[R]{{{lead-author-surname}} {\it et al}\ }
\fancyhead[C]{}

% Anti-widow / orphan
\widowpenalty=10000
\clubpenalty=10000
\displaywidowpenalty=10000
\raggedbottom

% Patch \articletype : replace the hardcoded "Journal Name" banner.
\makeatletter
\renewcommand{\articletype}[1]{%
   {\vspace*{-8mm}\noindent \Large \sf {{journal}}}%
   \vspace*{8mm} \noindent\reversemarginpar
   \marginpar{\vspace{-3mm} {\color{gray}\hrule} \ \\ Crossmark\\
       {\color{gray}\hrule} \ \\
       \tiny {\sf RECEIVED} {\small \\ dd Month yyyy}\\ \\
       {\sf REVISED} {\small \\ dd Month yyyy}}%
   {\scriptsize \sf{\bfseries \MakeUppercase{#1}}}}
\makeatother

% Theorem environments
\theoremstyle{plain}
\newtheorem{theorem}{Theorem}
\newtheorem{proposition}[theorem]{Proposition}
\newtheorem{lemma}[theorem]{Lemma}
\newtheorem{corollary}[theorem]{Corollary}
\theoremstyle{definition}
\newtheorem{definition}[theorem]{Definition}
\theoremstyle{remark}
\newtheorem{remark}[theorem]{Remark}
\newtheorem{conjecture}[theorem]{Conjecture}

% Bold-math shortcuts (delete what you don't need)
\newcommand{\E}{\mathbb{E}}
\newcommand{\Var}{\mathrm{Var}}
\newcommand{\by}{\mathbf{y}}
\newcommand{\bX}{\mathbf{X}}
\newcommand{\bphi}{\boldsymbol{\phi}}

% Editorial-placeholder marker used while drafting (renders less alarmingly
% than \textbf{[TODO]} ; reviewers will read it as a release-candidate
% placeholder rather than a substantive gap).  Remove all uses before v1.0.
\newcommand{\Todo}[1]{\textbf{[To confirm: #1]}}

% =========================================================================
\begin{document}

\articletype{Paper}     % or "Letter", "Topical Review"

\title{{{paper-title}}}

\author{{{lead-author-name}}$^{1,*}$\orcid{{{lead-author-orcid-id}}},
        {{co-author-name}}$^{2}$\orcid{{{co-author-orcid-id}}}}

\affil{$^{1}$ {{lead-author-affiliation}}
       \quad$^{2}$ {{co-author-affiliation}}\\[0.2em]
       $^{*}$ Author to whom any correspondence should be addressed.}

\email{{{lead-author-email}}}

\keywords{{{keyword1}}; {{keyword2}}; {{keyword3}}; {{keyword4}}; {{keyword5}}}

\begin{abstract}
\Todo{Write the abstract.  Target $\sim 250$ words.  Open with the
problem in one sentence, state the contribution in the second, give
the headline number (effect size + significance) in the third, and
close with the operational reading.}
\end{abstract}

% =========================================================================
\section{Introduction}
\label{sec:intro}

\Todo{Write Introduction.  Three paragraphs.  (1) The problem and why
existing approaches are insufficient.  (2) Our contribution at the
sentence-level.  (3) Roadmap of the paper.}

% =========================================================================
\section{Setup and notation}
\label{sec:setup}

\Todo{Formal setup.  Define the objects, the loss, the protocol.}

% =========================================================================
\section{Main result}
\label{sec:main}

\Todo{State the main theorem / claim with a half-page proof sketch ;
defer the full proof to the SI if longer than two pages.}

% =========================================================================
\section{Empirical validation}
\label{sec:empirics}

\Todo{Multi-pillar evidence architecture works well : (i) primary
predictive validation, (ii) robustness pillar, (iii) encoder
discrimination, (iv) falsifiability test against a topology-matched
null.  Convergence of independent statistics is the warrant.}

% =========================================================================
\section{Related work}
\label{sec:related}

\Todo{Position the contribution at the empty intersection of two
existing literatures.  One paragraph per closest adjacent result with
explicit "what we add" sentence.}

% =========================================================================
\section{Discussion}
\label{sec:discussion}

\Todo{Operational reading + connections to adjacent fields.}

% =========================================================================
\section{Limitations and future work}
\label{sec:limits}

\Todo{Be explicit about scope.  Address the obvious hostile-reviewer
objections (graph-construction sensitivity, sample size, benchmark
restrictions) with concrete theorem-level defenses where possible.}

% =========================================================================
\section{Methods}
\label{sec:methods}

\Todo{Reproducibility-grade Methods.  Data, encoders, splits,
statistical procedures, software stack with version pins.}

% =========================================================================
% Back matter -- iopjournal native commands.  Each takes its content as
% a single argument and emits the appropriate heading.  Under
% [anonymous] the contents (except \suppdata) are auto-redacted.
% =========================================================================

\ack{\Todo{Acknowledgements text.}}

\funding{\Todo{Funding sources and grant numbers.}}

\roles{\Todo{CRediT role assignments per author.}}

\data{All data sources and per-experiment derived outputs are
released under the MIT licence at
\url{https://github.com/cycling-data-lab/{{repo-name}}}.  A frozen
Zenodo snapshot of the v1.0 release tag is deposited at
\url{https://doi.org/10.5281/zenodo.XXXXXXX} \Todo{substitute the
issued Zenodo DOI after the first release}.}

\section*{Reproducibility}
Each table and figure is produced by a numbered script
\texttt{experiments/d<NN>\_<name>.py} that consumes a known input
cache, writes to a named output JSON, and pins \texttt{SEED=42}.
\Todo{Add experiment-specific reproducibility detail here.}

\section*{Competing interests}
The authors declare no competing financial or non-financial interests.

% =========================================================================
\section*{References}
\bibliography{references/references}

\end{document}
